\documentclass[11pt]{article}
\usepackage[margin=1in]{geometry}
\usepackage[T1]{fontenc}
\usepackage[utf8]{inputenc}
\usepackage{lmodern}
\usepackage{setspace}
\usepackage{microtype}
\usepackage[hidelinks]{hyperref}
\usepackage{parskip}
\usepackage{fancyhdr}

\setstretch{1.02}
\pagestyle{fancy}
\fancyhf{}
\lhead{Piter Garcia}
\rhead{EDU 498 Literacy Narrative}
\cfoot{\thepage}
\renewcommand{\headrulewidth}{0.3pt}

\title{Literacy Across Languages, Bodies, and Systems:\\A Personal and Community Literacy Narrative}
\author{Piter Garcia\\Literacy Learning as Social Practice, EDU 498}
\date{May 25, 2026}

\begin{document}
\maketitle
\thispagestyle{fancy}

\section*{Personal Literacy Autobiography}

My literacy story begins in the Dominican Republic, before I had the language I now use to explain myself. I first learned to read and write in Spanish through family expectations, school routines, copying, reading aloud, and the ordinary pressure to understand what adults wanted from me. At first, literacy looked simple from the outside: letters, notebooks, homework, pronunciation, and correctness. Underneath those school practices, though, literacy was also about belonging. It was about knowing how to sound capable, how to follow directions quickly enough, and how to avoid being misunderstood. As a child, I was often read through extremes: very capable in some moments, scattered or inconsistent in others. Only later did I begin to understand how neurodivergence, chronic illness, migration, grief, and cultural translation shaped the way I learned to read the world. EDU 498 has helped me name something I had already lived: literacy is not only an individual skill. It is a social practice shaped by language, culture, power, identity, and the tools available to us.

Moving to the United States at eighteen changed my relationship with literacy again. English was not only another language; it was a gatekeeping system. It shaped how professors heard me, how institutions sorted me, how forms expected me to explain myself, and how much energy it took to be seen as intelligent. I learned that academic English is not neutral. It rewards certain ways of organizing thought and often hides the labor of students who are translating across languages, cultures, bodies, and expectations. At the same time, my technical background in computer science, engineering, data science, and bioinformatics expanded my sense of literacy. I learned to read code, data, models, interfaces, scientific papers, medical language, and bureaucratic documents. I also learned that these literacies have consequences: they influence who receives support, who is believed, who is diagnosed, whose privacy is protected, and whose knowledge counts. Today, I view literacy as the ability to make meaning and claim agency inside social systems. For me, that includes Spanish and English, code and data, academic writing and family translation, medical self-advocacy and classroom teaching. It also includes literacies schools often miss: debugging, oral explanation, visual reasoning, peer support, technological navigation, and survival inside systems that were not designed for one's full complexity.

\section*{Community Exploration}

For the community portion of this narrative, I draw on a conversation with an older Dominican family member, whom I will refer to as a family elder for privacy. Their literacy journey began in a very different context from mine. Schooling was more limited, more teacher-directed, and more closely tied to discipline, handwriting, memorization, religious or moral instruction, and practical survival. Reading and writing mattered, but they were not always framed as personal expression or critical inquiry. Literacy was a tool for moving through everyday life: reading documents, understanding instructions, helping family members, managing work responsibilities, and later navigating migration, institutions, and technology.

When we talked about foundational literacy experiences, what stood out was how much literacy was connected to responsibility. For them, reading and writing were not only individual achievements; they were family and community resources. Someone who could read a letter, fill out a form, interpret official language, or help someone else understand a document held a kind of social power. This was especially true in contexts where migration, class, and language shaped access. Literacy could open doors, but it could also expose inequality. If the document was in the wrong language, if the institution assumed too much, or if the person asking for help felt shame, literacy became a reminder of who systems were designed to serve.

Technology also changed their literacy practices over time. They described a shift from paper-based literacy to phone-based literacy: text messages, WhatsApp, voice notes, online forms, video calls, digital banking, medical portals, and translation tools. These newer literacies are sometimes treated as simple or informal, but they are not. They require judgment, trust, interpretation, privacy awareness, and the ability to move between oral, written, visual, and digital modes. This connects directly to the EDU 498 emphasis on multiliteracies: people do not make meaning only through printed words. They make meaning across modes, communities, histories, and technologies.

Our literacy journeys are different, but they echo each other. I had access to more formal education, advanced technology, and academic pathways. I learned to read research articles, code, and data systems. My family elder learned literacy through necessity, migration, family responsibility, and practical problem-solving. Yet both stories involve translation, not only between Spanish and English, but between worlds. We both learned that literacy is tied to dignity. We both learned that institutions can make intelligent people feel inadequate when their language, accent, body, documentation, or mode of expression does not match what the system expects.

The comparison also helps me think about culturally and linguistically diverse students. A CLD student is not an empty container waiting for academic literacy. They may already be translating for family members, reading social cues across cultures, using digital tools to communicate, interpreting adult expectations, and navigating identities that school has not made space for. Similarly, a neurodivergent student may struggle with handwriting, timing, organization, or a multi-step direction while still showing powerful reasoning, pattern recognition, creativity, and insight. My family elder's story and my own story both challenge the idea that literacy development is a straight line. It is relational, uneven, embodied, multilingual, and shaped by power.

\section*{Reflection: Literacy as Social Practice}

Thinking across these two narratives pushes me to understand literacy as a social practice rather than a fixed individual skill. Literacy happens in relation to people, tools, histories, institutions, and identities. It is shaped by who is allowed to speak, whose language is treated as legitimate, whose writing is considered professional, whose questions are welcomed, and whose struggles are interpreted with generosity. Literacy is never just the ability to decode words. It is also the ability to participate, interpret power, advocate, create, and be recognized as a meaning-maker.

This course's focus on literacy learning as social practice helps me connect personal narrative to classroom design. If literacy is social, then teachers have to ask whose literacies are visible and whose are being ignored. If literacy is cultural, then students' languages, families, technologies, neurotypes, interests, and community knowledge are not distractions from learning; they are part of the learning. If literacy is critical, then students should not only consume texts and tools. They should also question how texts, technologies, policies, and institutions shape access, authorship, privacy, and belonging.

This matters for my future teaching, especially as a K-12 computer science teacher. In my Pine Brook IGNITE placement and Minecraft unit work, I have seen students demonstrate literacy through digital navigation, block coding, debugging, design choices, peer feedback, and oral explanation. A student may struggle to follow a click-path or written direction while still understanding the logic of a task. Another student may not write a polished paragraph but may explain a strategy clearly to a peer. If I only assess print-based products, I miss the literacy already happening in the room.

This is where my own experiences with neurodivergence, CLD identity, chronic illness, and multiple forms of academic and technical literacy become important. They remind me not to confuse access barriers with thinking barriers. A student who needs directions chunked, visuals repeated, language clarified, extra time, or another way to show learning is not less capable. They may simply be navigating a classroom system that has not yet made their literacy visible. Literacy instruction, then, should not be about forcing every student into one mode of expression. It should be about designing pathways where students can read, write, speak, listen, build, code, revise, question, and create in ways that maintain rigor while honoring complexity.

My community member's story also reminds me that literacy carries social consequences beyond school. The ability to read a form, question a policy, understand a diagnosis, interpret an algorithmic recommendation, or communicate across languages can affect health, work, family stability, and belonging. This is why literacy education cannot be separated from justice. When schools expand what counts as literacy, they expand who gets recognized as intelligent. When schools narrow literacy to one dominant language, one form of writing, or one pace of performance, they reproduce exclusion.

The lesson I take from these narratives is that literacy lives in the relationship between identity and participation. My literacy journey taught me to survive and translate across systems. My family elder's literacy journey taught me that reading and writing are community resources shaped by migration, responsibility, and adaptation. Together, they help me see literacy as a practice of dignity. In my classroom, I want students - especially neurodivergent students, CLD students, and students with multiple intersecting identities - to experience literacy as something they already possess and can continue to grow. I want them to leave with stronger skills, but also with a stronger sense that their languages, minds, bodies, and stories matter.

\end{document}
